\documentclass[11pt,a4paper]{article}

% ── FONTS ────────────────────────────────────────────────────
\usepackage{fontspec}
\setmainfont{Calibri}
\setsansfont{Calibri}

% ── LAYOUT ───────────────────────────────────────────────────
\usepackage[top=1.8cm,bottom=2cm,left=2cm,right=2cm]{geometry}
\usepackage{graphicx}
\usepackage{xcolor}
\usepackage{booktabs}
\usepackage{tabularx}
\usepackage{enumitem}
\usepackage{setspace}
\usepackage{fancyhdr}
\usepackage{titlesec}
\usepackage{hyperref}
\usepackage{microtype}
\usepackage{array}
\usepackage{colortbl}

\hypersetup{colorlinks=true,linkcolor=navyblue,urlcolor=accentteal,citecolor=navyblue}

% ── BRAND COLORS ────────────────────────────────────────────
\definecolor{navyblue}{HTML}{003566}
\definecolor{gold}{HTML}{C9A84C}
\definecolor{lightgray}{HTML}{F4F6FB}
\definecolor{darkgray}{HTML}{4A4A4A}
\definecolor{accentteal}{HTML}{0077B6}

% ── SECTION STYLING ─────────────────────────────────────────
\titleformat{\section}[block]
  {\normalfont\bfseries\fontsize{11}{13}\selectfont\color{navyblue}}
  {}{0em}{}
\titlespacing*{\section}{0pt}{10pt}{4pt}

\titleformat{\subsection}[block]
  {\normalfont\bfseries\fontsize{10}{12}\selectfont\color{accentteal}}
  {}{0em}{}
\titlespacing*{\subsection}{0pt}{6pt}{2pt}

% ── HEADER / FOOTER ─────────────────────────────────────────
\setlength{\headheight}{16pt}
\pagestyle{fancy}
\fancyhf{}
\fancyhead[L]{\footnotesize\color{navyblue}\textbf{AIGGPA Bhopal} \quad\textbar\quad
  Human Capital: Investment in Education, Skill Development and Employment}
\fancyhead[R]{\footnotesize\color{navyblue}April 2026}
\fancyfoot[C]{\footnotesize\color{navyblue}\thepage}
\renewcommand{\headrulewidth}{0.4pt}
\renewcommand{\headrule}{\color{navyblue}\hrule width\headwidth height\headrulewidth}
\renewcommand{\footrulewidth}{1pt}
\renewcommand{\footrule}{\color{gold}\hrule width\headwidth height 1pt}

% ── LISTS ───────────────────────────────────────────────────
\setlist[itemize]{leftmargin=1.2em, itemsep=1pt, topsep=2pt}

% ────────────────────────────────────────────────────────────
\begin{document}
\onehalfspacing

\begin{center}
  {\fontsize{12}{14}\selectfont\bfseries\color{navyblue}
  Youth Definitions, Indian \& MP Demography, and\\[2pt]
  Aspirations Towards Education, Skill \& Employment}\\[6pt]
  {\small\color{darkgray} Shreya \quad\textbar\quad AIGGPA Summer Intern 2026}
\end{center}

\vspace{0.2cm}

\section{1. Youth Definitions Across Agencies}

The term ``youth'' has no universal definition; it varies by agency and policy purpose:

\vspace{0.1cm}
\renewcommand{\arraystretch}{1.3}
\begin{table}[h!]
\centering
\footnotesize
\caption{\textbf{Youth Definitions -- National \& International Agencies}}
\vspace{0.1cm}
\begin{tabularx}{\textwidth}{>{\columncolor{lightgray}}l X >{\centering\arraybackslash}p{2.2cm}}
\toprule
\rowcolor{navyblue}
\textcolor{white}{\textbf{Agency / Framework}} & \textcolor{white}{\textbf{Definition of Youth}} & \textcolor{white}{\textbf{Age Group}} \\
\midrule
United Nations (UN)                     & Persons between 15 and 24 years (statistical purposes)         & 15--24 \\
\rowcolor{white}
WHO / UNICEF / UNFPA                    & ``Adolescents'' (10--19), ``Young People'' (10--24), ``Youth'' (15--24)  & 10--24 \\
Government of India (NYP 2014)          & National Youth Policy defines youth as persons aged 15--29     & 15--29 \\
\rowcolor{white}
Commonwealth Youth Programme            & Persons between 15 and 29 years                                & 15--29 \\
African Union                           & Persons between 15 and 35 years                                & 15--35 \\
\rowcolor{white}
\textbf{Generation Z (Gen Z)}            & Born between late 1990s and early 2010s ($\sim$1997--2012)     & $\sim$14--29 \\
\bottomrule
\end{tabularx}
\vspace{0.05cm}
\par\scriptsize Source: UN DESA; WHO; Ministry of Youth Affairs \& Sports, GOI (NYP 2014); MoSPI.
\end{table}

\textbf{Note on Gen Z:} India has the world's largest Gen Z cohort -- approx. \textbf{374--472 million} individuals, constituting $\sim$27\% of the population. This generation is digitally native, career-aspirational, and the primary target of both NEP~2020 and national skilling missions.

\vspace{0.15cm}
\section{2. Demography: India vs. Madhya Pradesh}

\renewcommand{\arraystretch}{1.25}
\begin{table}[h!]
\centering
\footnotesize
\caption{\textbf{Demographic Comparison -- India vs. Madhya Pradesh (2025 Projections)}}
\vspace{0.1cm}
\begin{tabularx}{\textwidth}{>{\columncolor{lightgray}}l >{\centering\arraybackslash}X >{\centering\arraybackslash}X}
\toprule
\rowcolor{navyblue}
\textcolor{white}{\textbf{Indicator}} & \textcolor{white}{\textbf{India}} & \textcolor{white}{\textbf{Madhya Pradesh}} \\
\midrule
Total Population (proj.)        & $\sim$144 crore                & $\sim$9 crore \\
\rowcolor{white}
Median Age                      & $\sim$28--29 years             & $\sim$25--26 years \\
Youth (15--29 yrs) share        & $\sim$27\%  ($\sim$370 Mn)     & $\sim$30\%  ($\sim$2.7 Cr) \\
\rowcolor{white}
Working Age (15--59) share      & $\sim$64\%                     & $\sim$59\% \\
Literacy Rate                   & 77.7\%                         & 69.3\% \\
\rowcolor{white}
Formal Vocational Training      & $\sim$2.4\% of 15+ pop.       & $\sim$1.2\% of 15+ pop. \\
\bottomrule
\end{tabularx}
\vspace{0.05cm}
\par\scriptsize Source: Census 2011 projections (2021 Census delayed); NSSO PLFS; Registrar General of India; ADB.
\end{table}

\vspace{0.1cm}
\section{3. Youth Aspirations -- Education, Skill \& Employment}

\begin{itemize}
  \item \textbf{Education:} Indian Gen Z shows a strong shift towards STEM, digital skills, and professional degrees. However, 45\% of graduates remain unemployable (India Skills Report 2025), indicating a \textit{quality gap} between aspiration and outcomes.
  \item \textbf{Skill Preferences:} AI/ML, Data Analytics, and Cloud Computing are the top three aspirational skill areas among youth aged 18--25. At the same time, 93\% of students desire internship opportunities as a bridge between academia and industry.
  \item \textbf{Employment:} While government jobs remain the top aspiration in states like MP (29 lakh registered ``aspirational youth'' on the state employment portal as of March 2025), urban youth increasingly aspire to startups, gig economy, and remote-work opportunities.
  \item \textbf{MP-specific:} MP's lower vocational training participation (1.2\% vs. 2.4\% national) means the state's youth aspirations often outpace available skill infrastructure. Initiatives like \textit{Mukhyamantri Seekho-Kamao Yojana} and the \textit{Global Skills Park, Bhopal} are designed to close this gap.
\end{itemize}

\vspace{0.25cm}
\noindent\rule{\textwidth}{0.5pt}
{\scriptsize\textbf{References:}
(1)~United Nations, DESA, \textit{Definition of Youth}, un.org/development/desa/youth.
(2)~Ministry of Youth Affairs \& Sports, \textit{National Youth Policy 2014}, yas.gov.in.
(3)~Ministry of Statistics \& Programme Implementation, \textit{Youth in India 2022}, mospi.gov.in.
(4)~India Skills Report 2025, CII--Wheebox, wheebox.com.
(5)~Asian Development Bank, \textit{MP Skills Development Project}, adb.org.
(6)~Registrar General of India, \textit{Population Projections 2011--2036}.
}

\end{document}
